% GENERATED FILE. Edit canonical lessons, then run npm run generate:books.
% canonical-source-hash: fnv1a64:52db9f79939754a1
% canonical-lessons: TE-C407-nirasana, TE-C407-udyamam, TE-C407-khandincu, TE-C407-dimandu, TE-C407-samskarana

\chapter{Protest and Reform}
\label{ch:te-protest-and-reform}
\hlchaptermodality{407}

\begin{chapteropening}
\textbf{By the end of this chapter:} \emph{I can talk about protests, movements, demands and reforms.}

Complete the last lesson of chapter 407: I can talk about protests, movements, demands and reforms.
\end{chapteropening}

\section[\texorpdfstring{\te{నిరసన} (nirasana)}{నిరసన (nirasana)}]{\te{నిరసన} (nirasana) — a protest}

\label{lesson:TE-C407-nirasana}



\begin{quote}
Before the new one: say the Telugu for an allegation, then the Telugu for to allege.
\end{quote}

\subsection*{You'll want to know: \te{నిరసన}}
\textbf{\te{నిరసన}} — \emph{nirasana} — \textquotedblleft{}a protest\textquotedblright{}.

\begin{quote}
\te{విద్యార్థులు} \te{కళాశాల} \te{ముందు} \te{నిరసన} \te{చేశారు}.
\end{quote}

\begin{quote}
\emph{vidyārthulu kaḷāśāla mundu nirasana cēśāru.}
\end{quote}

\begin{quote}
The students held a protest in front of the college.
\end{quote}

Plural: \te{నిరసనలు}, \emph{nirasanalu}.

\te{నిరసన} \te{తెలుపు}, to register a protest.

\begin{cousinweb}
\textbf{\te{నిరసన}} (\emph{nirasana}) is borrowed from Sanskrit \emph{nirasana}, \textquotedblleft{}casting out, driving away, rejecting\textquotedblright{} (Monier-Williams).
\end{cousinweb}

\begin{grammarlens}[title={What you've built}]
One more word for giving reasons and opinions.
\end{grammarlens}

\subsection*{Your turn}
\begin{itemize}
\raggedright
  \item \emph{Say it:} \emph{nirasana}
  \item \emph{Say it:} \emph{nirasana}, once more
  \item \emph{Say it:} \emph{vidyārthulu kaḷāśāla mundu nirasana cēśāru.} — The students held a protest in front of the college.
  \item \emph{Recall:} say the Telugu for an allegation, then the Telugu for to allege, then say \emph{nirasana} again
\end{itemize}

\begin{tcolorbox}[breakable,colback=teal!4,colframe=teal!35!black,title={Before you move on}]
What does \te{నిరసన} mean? (\textquotedblleft{}A protest\textquotedblright{}.) Say it once more.
\end{tcolorbox}
\section[\texorpdfstring{\te{ఉద్యమం} (udyamaṁ)}{ఉద్యమం (udyamaṁ)}]{\te{ఉద్యమం} (udyamaṁ) — a movement; a campaign for change}

\label{lesson:TE-C407-udyamam}



\begin{quote}
Before the new one: say the Telugu for to allege, then the Telugu for a protest.
\end{quote}

\subsection*{You'll want to know: \te{ఉద్యమం}}
\textbf{\te{ఉద్యమం}} — \emph{udyamaṁ} — \textquotedblleft{}a movement; a campaign for change\textquotedblright{}.

\begin{quote}
\te{ఆ} \te{గ్రామంలో} \te{నీళ్ళ} \te{కోసం} \te{ఒక} \te{ఉద్యమం} \te{మొదలైంది}.
\end{quote}

\begin{quote}
\emph{ā grāmaṁlō nīḷḷa kōsaṁ oka udyamaṁ modalaindi.}
\end{quote}

\begin{quote}
A movement for water began in that village.
\end{quote}

Plural: \te{ఉద్యమాలు}, \emph{udyamālu}.

A long campaign by many people. \te{ఉద్యమకారుడు}, an activist.

\begin{cousinweb}
\textbf{\te{ఉద్యమం}} (\emph{udyamaṁ}) is borrowed from Sanskrit \emph{udyama}, \textquotedblleft{}the act of raising\textquotedblright{}, and \textquotedblleft{}undertaking, striving after, zeal\textquotedblright{} (CDIAL 2048).
\end{cousinweb}

\begin{grammarlens}[title={What you've built}]
One more word for giving reasons and opinions.
\end{grammarlens}

\subsection*{Your turn}
\begin{itemize}
\raggedright
  \item \emph{Say it:} \emph{udyamaṁ}
  \item \emph{Say it:} \emph{udyamaṁ}, once more
  \item \emph{Say it:} \emph{ā grāmaṁlō nīḷḷa kōsaṁ oka udyamaṁ modalaindi.} — A movement for water began in that village.
  \item \emph{Recall:} say the Telugu for to allege, then the Telugu for a protest, then say \emph{udyamaṁ} again
\end{itemize}

\begin{tcolorbox}[breakable,colback=teal!4,colframe=teal!35!black,title={Before you move on}]
What does \te{ఉద్యమం} mean? (\textquotedblleft{}A movement; a campaign for change\textquotedblright{}.) Say it once more.
\end{tcolorbox}
\section[\texorpdfstring{\te{ఖండించు} (khaṇḍiñcu)}{ఖండించు (khaṇḍiñcu)}]{\te{ఖండించు} (khaṇḍiñcu) — to condemn; to denounce}

\label{lesson:TE-C407-khandincu}



\begin{quote}
Before the new one: say the Telugu for a protest, then the Telugu for a movement.
\end{quote}

\subsection*{You'll want to know: \te{ఖండించు}}
\textbf{\te{ఖండించు}} — \emph{khaṇḍiñcu} — \textquotedblleft{}to condemn; to denounce\textquotedblright{}.

\begin{quote}
\te{పెద్దలు} \te{ఆ} \te{పనిని} \te{గట్టిగా} \te{ఖండించారు}.
\end{quote}

\begin{quote}
\emph{peddalu ā panini gaṭṭigā khaṇḍiñcāru.}
\end{quote}

\begin{quote}
The elders strongly condemned that act.
\end{quote}

Past (I …): \te{ఖండించాను}, \emph{khaṇḍiñcānu}.

Also to refute: \te{ఆ} \te{వాదనను} \te{ఖండించు}, to refute that argument.

\begin{cousinweb}
\textbf{\te{ఖండించు}} (\emph{khaṇḍiñcu}) is built on Sanskrit \emph{khaṇḍana}, \textquotedblleft{}breaking\textquotedblright{} (CDIAL 3796), with the Telugu verb ending -\te{ఇంచు}.
\end{cousinweb}

\begin{grammarlens}[title={What you've built}]
One more word for giving reasons and opinions.
\end{grammarlens}

\subsection*{Your turn}
\begin{itemize}
\raggedright
  \item \emph{Say it:} \emph{khaṇḍiñcu}
  \item \emph{Say it:} \emph{khaṇḍiñcu}, once more
  \item \emph{Say it:} \emph{peddalu ā panini gaṭṭigā khaṇḍiñcāru.} — The elders strongly condemned that act.
  \item \emph{Recall:} say the Telugu for a protest, then the Telugu for a movement, then say \emph{khaṇḍiñcu} again
\end{itemize}

\begin{tcolorbox}[breakable,colback=teal!4,colframe=teal!35!black,title={Before you move on}]
What does \te{ఖండించు} mean? (\textquotedblleft{}To condemn; to denounce\textquotedblright{}.) Say it once more.
\end{tcolorbox}
\section[\texorpdfstring{\te{డిమాండు} (ḍimāṇḍu)}{డిమాండు (ḍimāṇḍu)}]{\te{డిమాండు} (ḍimāṇḍu) — a claim; a demand}

\label{lesson:TE-C407-dimandu}



\begin{quote}
Before the new one: say the Telugu for a movement, then the Telugu for to condemn.
\end{quote}

\subsection*{You'll want to know: \te{డిమాండు}}
\textbf{\te{డిమాండు}} — \emph{ḍimāṇḍu} — \textquotedblleft{}a claim; a demand\textquotedblright{}.

\begin{quote}
\te{ఉద్యోగుల} \te{డిమాండు} \te{ఎక్కువ} \te{జీతం}.
\end{quote}

\begin{quote}
\emph{udyōgula ḍimāṇḍu ekkuva jītaṁ.}
\end{quote}

\begin{quote}
The employees’ demand is higher pay.
\end{quote}

Plural: \te{డిమాండ్లు}, \emph{ḍimāṇḍlu}.

Like \te{బస్సు}, it ends in -\te{ఉ}. \te{డిమాండు} \te{చేయు}, to demand.

\begin{cousinweb}
\textbf{\te{డిమాండు}} (\emph{ḍimāṇḍu}) is a loan from English \emph{demand}, with the final -\te{ఉ} that English loans often take.
\end{cousinweb}

\begin{grammarlens}[title={What you've built}]
One more word for giving reasons and opinions.
\end{grammarlens}

\subsection*{Your turn}
\begin{itemize}
\raggedright
  \item \emph{Say it:} \emph{ḍimāṇḍu}
  \item \emph{Say it:} \emph{ḍimāṇḍu}, once more
  \item \emph{Say it:} \emph{udyōgula ḍimāṇḍu ekkuva jītaṁ.} — The employees’ demand is higher pay.
  \item \emph{Recall:} say the Telugu for a movement, then the Telugu for to condemn, then say \emph{ḍimāṇḍu} again
\end{itemize}

\begin{tcolorbox}[breakable,colback=teal!4,colframe=teal!35!black,title={Before you move on}]
What does \te{డిమాండు} mean? (\textquotedblleft{}A claim; a demand\textquotedblright{}.) Say it once more.
\end{tcolorbox}
\section[\texorpdfstring{\te{సంస్కరణ} (saṁskaraṇa)}{సంస్కరణ (saṁskaraṇa)}]{\te{సంస్కరణ} (saṁskaraṇa) — a reform}

\label{lesson:TE-C407-samskarana}



\begin{quote}
Before the new one: say the Telugu for to condemn, then the Telugu for a claim.
\end{quote}

\subsection*{You'll want to know: \te{సంస్కరణ}}
\textbf{\te{సంస్కరణ}} — \emph{saṁskaraṇa} — \textquotedblleft{}a reform\textquotedblright{}.

\begin{quote}
\te{విద్యా} \te{రంగంలో} \te{కొత్త} \te{సంస్కరణలు} \te{వచ్చాయి}.
\end{quote}

\begin{quote}
\emph{vidyā raṅgaṁlō kotta saṁskaraṇalu vaccāyi.}
\end{quote}

\begin{quote}
New reforms have come in education.
\end{quote}

Plural: \te{సంస్కరణలు}, \emph{saṁskaraṇalu}.

Mostly in the plural, \te{సంస్కరణలు}.

\begin{cousinweb}
\textbf{\te{సంస్కరణ}} (\emph{saṁskaraṇa}) is borrowed from Sanskrit \emph{saṁskaraṇa}, \textquotedblleft{}the act of putting together, preparing\textquotedblright{} (Monier-Williams); Telugu uses it for putting a system right.
\end{cousinweb}

\begin{grammarlens}[title={What you've built}]
One more word for giving reasons and opinions.
\end{grammarlens}

\subsection*{Your turn}
\begin{itemize}
\raggedright
  \item \emph{Say it:} \emph{saṁskaraṇa}
  \item \emph{Say it:} \emph{saṁskaraṇa}, once more
  \item \emph{Say it:} \emph{vidyā raṅgaṁlō kotta saṁskaraṇalu vaccāyi.} — New reforms have come in education.
  \item \emph{Recall:} say the Telugu for to condemn, then the Telugu for a claim, then say \emph{saṁskaraṇa} again
\end{itemize}

\begin{tcolorbox}[breakable,colback=teal!4,colframe=teal!35!black,title={Before you move on}]
What does \te{సంస్కరణ} mean? (\textquotedblleft{}A reform\textquotedblright{}.) Say it once more.
\end{tcolorbox}
